\section{Evaluation protocol}

The protocol is designed to turn the compiler into a measuring instrument, not
to imply that implementation completeness is a scientific result.

We preregister three questions. RQ1 asks how often source-grounded candidates
survive the bounded compiler and which constructs are refused. RQ2 asks whether
machine checks and independent execution agree on the supported subset. RQ3
asks whether an artifact-free source signal predicts outcome equivalence after
conditioning on model, system, and run. The corresponding hypotheses are
directional but not assumed to hold: H1 predicts that unsupported and malformed
constructs are concentrated in explicit refusal/error classes; H2 predicts that
independent-engine disagreements are explainable by a finite set of lowering
or semantic-boundary causes; and H3 predicts positive association between AFS
and OE under the preregistered controls. A null, weak, underpowered, or invalid
outcome is a valid result and triggers the stop rules in Section 7.

\subsection{Corpora and units}

The primary unit is a source-bearing rule, not a generated sentence or model
call. A licensed, stratified frame must contain document spans, rule labels,
exception and threshold annotations, and dependency edges. Dutch regulatory
text is the planned natural-language anchor; ContractNLI is a planned
cross-domain adapter. Privacy-policy runs are useful operational stress tests
but cannot stand in for the licensed annotation frame.

\subsection{Measurements and controls}

AFS is measured with two independent annotators, adjudication, weighted
precision/recall, and agreement on presence, value, exception, and evidence
span. OE uses frozen source models, generated models, and a pinned independent
DMN engine; exact outcomes and disagreement classes are retained. Controls
include source-only, generated-only, and executable-oracle conditions. Prompt,
model, temperature, retry, and batch settings are pinned. Exploratory pilots
are labeled and excluded from headline intervals.

For a binomial proportion, we report Wilson intervals and the realized
denominator. Bootstrap intervals are paired by document for agreement and
latency summaries. We preregister handling of refusals, missing evidence,
unresolved references, and ties before unsealing labels. Multiple comparisons
are controlled within a declared family; a stopped or underpowered study is
reported as such.

We report executable yield and conditional quality separately. If $U$ is the
retained rule universe and $E$ the rules that lower and execute, executable
yield is $EY=|E|/|U|$. Quality is reported conditional on $E$ (for example,
outcome agreement or source recall); refusals and invalid cases remain in the
yield denominator. This avoids turning a selective compiler into an apparently
accurate system by silently dropping hard cases.

\subsection{Reproducibility contract}

Each result ships a manifest containing source and artifact hashes, code commit,
environment lock, model/provider configuration, random seeds, exclusion log,
and validator output. A replay is valid only if schema versions and evidence
packet semantics match. The release artifact distinguishes code that is
executable locally from observations that require licensed data, provider
access, GPU time, or human adjudication.
